% Fira Code 6.2, embedded by pdfTeX with the stock TeXBase1 encoding.
\pdfmapline{+FiraCode-Regular FiraCode-Regular "TeXBase1Encoding ReEncodeFont" <8r.enc <../../latex-commons/fonts/firacode/FiraCode-Regular.ttf}
\pdfmapline{+FiraCode-Bold FiraCode-Bold "TeXBase1Encoding ReEncodeFont" <8r.enc <../../latex-commons/fonts/firacode/FiraCode-Bold.ttf}
\DeclareFontFamily{T1}{answerfiracode}{}
\DeclareFontShape{T1}{answerfiracode}{m}{n}{<-> ../../latex-commons/fonts/firacode/FiraCode-Regular}{}
\DeclareFontShape{T1}{answerfiracode}{b}{n}{<-> ../../latex-commons/fonts/firacode/FiraCode-Bold}{}
\DeclareFontShape{T1}{answerfiracode}{bx}{n}{<-> ../../latex-commons/fonts/firacode/FiraCode-Bold}{}

% Keep TeX's typewriter character handling within this listing.
\newcommand{\answercodefont}{%
  \renewcommand{\ttdefault}{answerfiracode}%
  \fontencoding{T1}\fontshape{n}\ttfamily\selectfont
  % TeXBase1 places the straight quote and backtick in slots 31 and 30.
  \renewcommand{\textquotesingle}{\char31\relax}%
  \renewcommand{\textasciigrave}{\char30\relax}%
}
\DeclareTextFontCommand{\answertexttt}{\answercodefont}

% Fira Code has no italic face. Shear the styled token while retaining its width.
\newsavebox{\answercodeslantbox}
\newcommand{\answercodeslanttoken}[1]{%
  \sbox{\answercodeslantbox}{#1}%
  \hbox to \wd\answercodeslantbox{%
    \pdfsave\pdfsetmatrix{1 0 .2 1}%
    \rlap{\usebox{\answercodeslantbox}}%
    \pdfrestore\hss
  }%
}

% Catppuccin Latte token and bracket colors from its VS Code theme.
\definecolor{AnswerCodeBase}{HTML}{EFF1F5}
\definecolor{AnswerCodeBorder}{HTML}{ACB0BE}
\definecolor{AnswerCodeText}{HTML}{4C4F69}
\definecolor{AnswerCodeComment}{HTML}{7C7F93}
\definecolor{AnswerCodeKeyword}{HTML}{8839EF}
\definecolor{AnswerCodeNumber}{HTML}{FE640B}
\definecolor{AnswerCodeOperator}{HTML}{179299}
\definecolor{AnswerCodeFunction}{HTML}{1E66F5}
\definecolor{AnswerCodeParameter}{HTML}{E64553}
\definecolor{AnswerCodeString}{HTML}{40A02B}
\definecolor{AnswerCodeBracketOne}{HTML}{D20F39}
\definecolor{AnswerCodeBracketThree}{HTML}{DF8E1D}
\definecolor{AnswerCodeBracketFive}{HTML}{209FB5}

\makeatletter
\newcommand{\answercodeslant}{\let\lst@alloverstyle\answercodeslanttoken}

% Inside a list, start at the label edge. Elsewhere use the text margin.
\newcommand{\answercodeleftskip}{\ifdim\@totalleftmargin>\z@ -\leftmargin\else 0pt\fi}

\newif\ifanswercodehighlight
\newcount\answercodebracketdepth
\newcommand{\answercodebracketcolor}{%
  \ifcase\answercodebracketdepth
  AnswerCodeBracketOne\or AnswerCodeNumber\or AnswerCodeBracketThree%
  \or AnswerCodeString\or AnswerCodeBracketFive\or AnswerCodeKeyword%
  \else AnswerCodeKeyword\fi
}
\newcommand{\answercodeopen}[1]{%
  \textcolor{\answercodebracketcolor}{#1}%
  \global\advance\answercodebracketdepth by 1\relax
  \ifnum\answercodebracketdepth=6 \global\answercodebracketdepth=0\fi
}
\newcommand{\answercodeclose}[1]{%
  \ifnum\answercodebracketdepth=0 \global\answercodebracketdepth=6\fi
  \global\advance\answercodebracketdepth by -1\relax
  \textcolor{\answercodebracketcolor}{#1}%
}

% listings' line wrapping replaces its normal closing-parenthesis handler.
\newcommand{\answercoderightparen}{%
  \lst@Def{`)}{%
    \lst@ifmode
      \lst@breakProcessOther)%
    \else
      \lst@XPrintToken
      \let\lst@scanmode\lst@scan@m
      \lst@token{{\answercodeclose{)}}}\lst@length\@ne
      \lst@XPrintToken
    \fi
  }%
}

% `listings` sends numbers and punctuation together to OutputOther. The
% literal syntax below separates operators and brackets before number styling.
\ExplSyntaxOn
\cs_new_protected:Npn \answercodestyleother
  {
    \tl_set:No \l_tmpa_tl { \the\lst@token }
    \exp_args:NnV \regex_match:nnTF
      { \A (?: [0-9]+ (?: \.[0-9]* )? | \.[0-9]+ ) \Z } \l_tmpa_tl
      { \def\lst@thestyle{\color{AnswerCodeNumber}} }
      { }
  }
\ExplSyntaxOff
\lst@AddToHook{OutputOther}{%
  \ifanswercodehighlight
    \lst@ifmode\else\answercodestyleother\fi
  \fi
}
\makeatother

% Paste code inline. Use a fixed-height tcolorbox without `fit` around it.
% The second optional argument accepts extra listings keys for semantic names.
\NewTCBListing{codeblock}{O{Python} O{}}{
  listing only,
  left skip=\answercodeleftskip,
  colback=AnswerCodeBase,
  colframe=AnswerCodeBorder,
  boxrule=0.4pt, arc=3pt,
  left=6pt, right=6pt, top=4pt, bottom=4pt,
  before skip=5pt, after skip=5pt,
  listing options={
    language=#1,
    basicstyle={\small\answercodefont\color{AnswerCodeText}%
        \answercodehighlighttrue\global\answercodebracketdepth=0\relax},
    keywordstyle=\color{AnswerCodeKeyword},
    commentstyle=\color{AnswerCodeComment}\answercodeslant,
    stringstyle=\color{AnswerCodeString},
    emphstyle={[1]\color{AnswerCodeParameter}\answercodeslant},
    emphstyle={[2]\color{AnswerCodeFunction}\answercodeslant},
    columns=fullflexible, keepspaces=true,
    upquote=true, showstringspaces=false, breaklines=true,
    xleftmargin=0pt, frame=none,
    backgroundcolor=\color{AnswerCodeBase},
    aboveskip=0pt, belowskip=0pt,
    MoreSelectCharTable=\answercoderightparen,
    literate=*
        {(}{{\answercodeopen{(}}}1
        {[}{{\answercodeopen{[}}}1 {]}{{\answercodeclose{]}}}1
        {\{}{{\answercodeopen{\textbraceleft}}}1
        {\}}{{\answercodeclose{\textbraceright}}}1
        {=}{{\textcolor{AnswerCodeOperator}{=}}}1
        {+}{{\textcolor{AnswerCodeOperator}{+}}}1
        {-}{{\textcolor{AnswerCodeOperator}{-}}}1
        {*}{{\textcolor{AnswerCodeOperator}{*}}}1
        {/}{{\textcolor{AnswerCodeOperator}{/}}}1
        {<}{{\textcolor{AnswerCodeOperator}{\textless}}}1
        {>}{{\textcolor{AnswerCodeOperator}{\textgreater}}}1
        {:}{{\textcolor{AnswerCodeComment}{:}}}1
        {,}{{\textcolor{AnswerCodeComment}{,}}}1,
    #2
  }
}
